%% Section 4: Discussion and Limitations
\section{Discussion and Limitations}
\label{sec:discussion}

The UML-driven framework presented in \Cref{sec:modeling} constitutes, to our knowledge, the first comprehensive formal modeling of a mobile AR-enhanced vocabulary learning system that integrates OCR scanning, intelligent tutoring, augmented reality overlays, and gamification within a unified pedagogical architecture. The fifteen classes organized across five functional groups---Core Actors, Learning Content, Student Modeling, AI \& AR, and Gamification---provide a clear separation of concerns in which each class bears a well-defined pedagogical responsibility. This structural clarity is not merely an engineering convenience; it is a deliberate design strategy that ensures every component can be traced back to a specific educational requirement. The \texttt{StudentModel} implements Bayesian Knowledge Tracing~\cite{Corbett1994}, the \texttt{TutoringEngine} operationalizes Bloom's taxonomy~\cite{Bloom1956,Krathwohl2002}, the \texttt{GamificationEngine} instantiates Keller's ARCS model~\cite{Keller2010}, and the \texttt{ARRenderer} embodies the affordances of situated, context-aware AR learning~\cite{Azuma1997,Akcayir2017}. The four complementary UML diagrams---class, use case, activity, and sequence---collectively bridge the gap between abstract pedagogical intent and concrete technical specification, providing a reusable blueprint that can guide both implementation and future research.

The proposed framework addresses three core challenges identified in the literature. \emph{First}, the vocabulary acquisition challenge is tackled through contextual, adaptive explanations that preserve the relationship between words and the passages in which they appear~\cite{Hulstijn2001,Nation2013}. By coupling OCR-based text extraction with ITS-generated explanations calibrated to the learner's mastery level and Bloom's taxonomy stage, AR-DeC ensures that learners receive definitions, examples, and related concepts that are neither too simple (causing boredom) nor too complex (causing frustration)~\cite{Laufer2010,Schmitt2008}. \emph{Second}, the engagement challenge is addressed through a gamification engine aligned with the ARCS motivational model~\cite{Keller2010}: AR visuals capture \emph{Attention}, contextual scanning ensures \emph{Relevance}, adaptive difficulty builds \emph{Confidence}, and points, badges, and streaks deliver \emph{Satisfaction}~\cite{Deterding2011,Hamari2014,Sailer2017}. The quiz component, with its timed questions and combo multipliers, transforms vocabulary review into an intrinsically motivating activity. \emph{Third}, the personalization challenge is resolved through Bayesian Knowledge Tracing, which maintains per-word mastery estimates ($P(L_n)$) updated after every interaction~\cite{Corbett1994}. The four BKT parameters---prior knowledge, transition probability, slip probability, and guess probability---enable the system to distinguish between genuine mastery and lucky guessing, providing the tutoring engine with the fine-grained knowledge state information it needs to select appropriate content~\cite{VanLehn2011}.

A distinctive contribution of the AR-DeC framework lies in how UML serves as a bridge between pedagogical design and technical implementation. Each UML diagram encodes specific pedagogical decisions that might otherwise remain implicit in the minds of designers. For instance, the tiered help-seeking mechanism---modeled as an \texttt{<<extend>>} relationship in the use case diagram and as an \texttt{alt} combined fragment in the sequence diagram---reflects a deliberate pedagogical choice: the conversational agent handles quick factual queries, the intelligent tutor provides deeper adaptive explanations, and the human teacher is reserved for genuinely challenging cases. This escalation hierarchy is not arbitrary; it is grounded in research on scaffolding and zone of proximal development, ensuring that learners receive the minimum necessary support to progress independently~\cite{VanLehn2011}. Similarly, the activity diagram's fork/join bars encode the parallel execution of mastery updates and gamification processing, a design decision that minimizes latency and ensures that the learner receives immediate motivational feedback. Our approach shares methodological foundations with the UML-driven pedagogical modeling proposed by Ouariach et al.~\cite{Ouariach2023,Ouariach2025,Ouariach2026} for collaborative e-learning systems, but extends it in three directions: (1)~the integration of on-device OCR as a content acquisition mechanism, (2)~the incorporation of AR as a presentation modality, and (3)~the embedding of BKT-based adaptive mastery tracking within the class structure. Whereas Ouariach et al. focused on collaborative learning workflows among groups of learners, AR-DeC centers on the individual learner's interaction with physical textbooks, augmented by intelligent, immersive, and gamified digital support.

Despite these contributions, several limitations must be acknowledged. \emph{First}, the framework presented here is a conceptual model that has not yet been implemented as a functional prototype or evaluated empirically with learners; consequently, the pedagogical effectiveness of the proposed architecture remains to be validated through controlled studies. \emph{Second}, OCR accuracy is inherently dependent on text quality---factors such as font size, print contrast, page curvature, and ambient lighting can significantly affect extraction reliability, and the current model does not specify detailed error-handling or confidence-thresholding mechanisms for degraded OCR conditions. \emph{Third}, the BKT parameters (prior knowledge, transition, slip, and guess probabilities) require calibration on real learner data; the default values commonly used in the literature may not generalize to all vocabulary domains or learner populations~\cite{Corbett1994}. \emph{Fourth}, the AR experience is inherently device-dependent: camera quality, processing power, and AR framework support vary across smartphones, potentially creating equity concerns if learners use devices of substantially different capabilities~\cite{Garzon2019}. \emph{Fifth}, while the UML diagrams specify \emph{what} the system should do and \emph{how} components interact, they deliberately abstract away certain implementation details---for instance, the specific AI model used for explanation generation, the rendering pipeline for AR overlays, and the backend synchronization protocol---which must be resolved during the engineering phase.

These limitations suggest several directions for future work. The most immediate next step is the implementation of a functional prototype based on the UML specifications, followed by a Design-Based Research (DBR) evaluation involving iterative design-evaluate-refine cycles with real learners in authentic educational settings~\cite{Crompton2017}. Empirical studies should assess the impact of AR-DeC on vocabulary acquisition, learner engagement, and knowledge retention, comparing the integrated system against control conditions (e.g., AR without ITS, ITS without gamification) to isolate the contribution of each component. Additionally, extending the framework to support multi-language vocabulary acquisition, collaborative learning features (e.g., peer quizzing, shared vocabulary lists), and teacher-authored AR content would broaden its applicability across educational contexts and learner populations. Finally, investigating the calibration of BKT parameters through machine learning on interaction logs would improve the accuracy of mastery estimates and, consequently, the quality of adaptive content selection.
